\section{Introduction}\label{sec:introduction}

For \(E\subset\R^d\), its Euclidean distance set is
\[
 \Delta(E)=\{|x-y|:x,y\in E\}.
\]
The Falconer distance conjecture asks whether a compact set of Hausdorff
dimension greater than half the ambient dimension must determine a set
of distances of positive length. We prove this assertion.

\begin{theorem}\label{thm:main}
Let \(d\ge2\) be an integer and let \(E\subset\R^d\) be compact. If
\(\dim_H E>d/2\), then
\[
 \cL^1(\Delta(E))>0.
\]
\end{theorem}

Theorem~\ref{thm:main} imposes no regularity beyond the strict Hausdorff
dimension assumption. In particular, it does not require equality of
Hausdorff and packing dimensions, an Ahlfors regular measure, product
structure, or Fourier decay. The assertion concerns positive Lebesgue
measure; a distance set of Hausdorff dimension one alone would not
suffice. The strict threshold is retained, and no endpoint assertion is made.

